\section{Conclusion}

This study shows that conclusions about common recognition heads and losses for palmprint embeddings depend strongly on the evaluation protocol. Under a seen-identity diagnostic setting, BNNeck+ArcFace (M6) can appear highly favorable, but under held-out palm-class Tongji cross-session evaluation it does not provide a consistent improvement over a cross-entropy plus supervised-contrastive baseline (M1). The strict protocol also reveals session-direction asymmetry and differences between Rank-1 identification and low-FAR verification behavior. 

The strict component ablation further shows that M4 (BNNeck + CE) has the highest observed mean on the selected Rank-1 and low-FAR TAR criteria among the tested learned variants, indicating that while BNNeck can shape useful embedding geometry under some session directions, ArcFace-based margins do not provide a reliable additional gain under the strict palm-class-disjoint protocol.

The corrected IITD palm-class-disjoint within-session secondary validation supports this conclusion, showing a near-tie between M1 and M6 rather than a clear M6 advantage. The fixed Gabor texture reference baseline further contextualizes these results, showing that classical texture descriptors can remain competitive for nearest-neighbor Rank-1 identification while performing much worse than learned embeddings in the strict low-FAR verification regime.

These results do not imply that BNNeck, ArcFace, or supervised contrastive learning are ineffective for palmprints; rather, they show that their apparent value is conditional on split design, session direction, and threshold policy. The main contribution is an audited, reproducible evaluation framework and a calibrated empirical finding: component-level palmprint claims should be tied to held-out class protocols and conservative biometric metrics.
